\section{Introduction}\label{sec:intro}

Software performance is a first-order quality attribute on mobile platforms, where the largest available headroom lies at the \emph{software} level: empirical studies consistently show that library choice, data-structure choice, and algorithm choice routinely deliver multi-fold improvements in latency and energy~\cite{patterns,reduce_bytecode_energy,marco_refactorings_wild,empirical_rua}. Among classical software-level optimisations, \emph{memoization}~\cite{michie1968memo}---caching the result of a pure method under its argument tuple---is one of the best understood, with empirical evidence of order-of-magnitude speedups on recursion-heavy workloads whose naive form visits overlapping subproblems~\cite{hall_memoization,costa_memoization,immutator}.

Despite this pedigree, applying memoization to an existing \jvm or Android codebase correctly remains a non-trivial refactoring. A correct memoized implementation must simultaneously address at least six concerns: (\circnum{1})~cache-field placement, (\circnum{2})~cache-key construction from argument tuples, (\circnum{3})~null handling, (\circnum{4})~per-row or bulk invalidation on mutation, (\circnum{5})~thread-safety, and (\circnum{6})~safe interaction with downstream code transformations---compile-time optimisers, build-time shrinkers (R8), and runtime class rewriters---that may rename, inline, or eliminate the synthetic fields implementing the cache. The recent iMMutator study~\cite{immutator} provides direct empirical evidence that developers trip over this correctness surface: hand-written memoization patches recurrently exhibit broken invalidation, thread-unsafe access, and interaction bugs with toolchain transformations. The study's conclusion---that \emph{detection} can be automated but \emph{application} remains a developer bottleneck---is the direct motivation for this work.

\paragraph{Existing caches do not close the gap.} Guava~\cite{guava} and Caffeine~\cite{caffeine} are server-side libraries whose builder-centric APIs require wrapper code at every call site. Spring Cache~\cite{spring_cache} has the right annotation ergonomics but is inseparable from the Spring container and absent on Android. Function-level wrapper libraries replace one refactoring with another. None of them let an automated tool make an existing method memoized \emph{in place}, without touching its call sites, its enclosing type, or the framework around it.

\paragraph{Contributions.} We present \textbf{MemoizeLib}, a memoization library whose design closes the detection--application gap. Its contributions are: (\circnum{1})~a three-annotation API---\code{@Memoize}, \code{@CacheInvalidate}, \code{@InvalidateCacheEntry}---that collapses the six-concern surface behind a single declarative boundary; (\circnum{2})~a cross-platform build-time bytecode transformation that rewrites annotated methods in place on both Android (via \agp) and the plain \jvm (via a standalone transformer), sharing a zero-dependency runtime module; (\circnum{3})~a runtime offering three eviction policies (\lru, \fifo, \lfu) plus an unbounded mode, with \ttl, auto-monitor, and embedded observability whose disabled-state overhead reduces to a single volatile read (derived in Section~\ref{sec:tool}); and (\circnum{4})~an open Jetpack Benchmark harness over ten overlapping-subproblem algorithms released as a reproducible research artefact for downstream detectors and refactoring tools.
